% ================================================================
%  Chapter 5 — Detailed Design (KANDA)
%  Knowledge-driven Autonomous Navigation and Decision-making Agent
% ================================================================
\chapter{Detailed Design}
\label{ch:detail}

This chapter presents the detailed design of \textbf{Knowledge-driven Autonomous Navigation and Decision-making Agent (KANDA)} at module and component level. It builds on the high-level architecture and data flows in Chapter~\ref{ch:hld}: overview diagrams, use cases, deployment components, a structure chart for the Python orchestrator, an activity view of one control cycle, and descriptions of each software module together with the embodiment firmware running on the ESP32.

%-------------------------------------------------------------------
\section{High-Level Design Overview}
\label{sec:hld_overview}

The deployment comprises a Raspberry Pi running the Python package tree \texttt{ai\_layer/}, an ESP32 running Arduino-compatible firmware under \texttt{firmware/phase2\_obstacle\_avoidance/}, a duplex UART link, and HTTPS access to Google Gemini (default model \texttt{gemini-2.5-flash}). Figure~\ref{fig:hld_overview} summarises the principal deployable blocks.

\begin{figure}[htbp!]
\centering
\resizebox{0.9\textwidth}{!}{%
\begin{tikzpicture}[
  block/.style={draw, thick, rounded corners=4pt, minimum width=3.6cm, minimum height=1.15cm, align=center, font=\small},
  arr/.style={-{Stealth[length=2.5mm]}, thick},
  darr/.style={-{Stealth[length=2.5mm]}, thick, densely dashed},
  lbl/.style={font=\scriptsize, fill=white, inner sep=2pt},
  node distance=1.1cm
]

\node[block, fill=blue!10] (op) {\textbf{Operator}\\{\footnotesize supervision}};

\node[block, fill=cyan!10, right=1.8cm of op] (pi) {\textbf{Raspberry Pi}\\{\footnotesize \texttt{ai\_layer/} (Python)}};

\node[block, fill=orange!12, right=1.8cm of pi] (esp) {\textbf{ESP32}\\{\footnotesize firmware executor}};

\node[block, fill=purple!10, below=1.45cm of pi] (gem) {\textbf{Gemini API}\\{\footnotesize HTTPS / JSON}};

\node[block, fill=gray!14, right=1.5cm of esp] (phys) {\textbf{Physical stack}\\{\footnotesize sensors, driver, motors, OLED}};

\draw[arr] (op) -- node[lbl, above] {start/stop} (pi);
\draw[arr] (pi) -- node[lbl, above] {UART JSON} (esp);
\draw[arr] (esp) -- node[lbl, above] {mounting} (phys);
\draw[arr] (pi.south) -- node[lbl, left] {inference} (gem.north);
\draw[arr] (gem.north) -- node[lbl, right] {completion} (pi.south);

\end{tikzpicture}%
}% end resizebox
\caption{High-level deployable blocks of KANDA}
\label{fig:hld_overview}
\end{figure}

Fig.~\ref{fig:hld_overview} shows how the operator interacts with the intelligence layer on the Pi, how the Pi exchanges newline-terminated JSON and telemetry with the ESP32 over UART, how the ESP32 connects to the physical stack, and how the Pi calls the Gemini API over the network.

%-------------------------------------------------------------------
\section{Low-Level Design Overview}
\label{sec:lld_overview}

The Python intelligence layer is organised into configuration, serial I/O, prompt construction, LLM access, validation, and a single orchestrator entry point. Figure~\ref{fig:lld_overview} maps files to these roles.

\begin{figure}[htbp!]
\centering
\resizebox{0.94\textwidth}{!}{%
\begin{tikzpicture}[
  mod/.style={draw, thick, rounded corners=3pt, minimum width=3.5cm, minimum height=0.85cm, align=center, font=\footnotesize},
  grp/.style={draw, thick, rounded corners=6pt, minimum width=8cm, minimum height=2.4cm},
  arr/.style={-{Stealth[length=2mm]}, thick},
  lbl/.style={font=\tiny, fill=white, inner sep=1pt},
  node distance=0.45cm
]

\node[mod, fill=blue!12] (main) {\textbf{main.py}\\{\scriptsize \texttt{run\_cycle()}, event loop}};
\node[mod, fill=green!10, below left=0.55cm and 0.2cm of main] (uart) {\textbf{uart\_bridge.py}\\{\scriptsize \texttt{UARTBridge}}};
\node[mod, fill=green!10, below=0.55cm of main] (ctx) {\textbf{context\_builder.py}\\{\scriptsize \texttt{build\_prompt()}}};
\node[mod, fill=green!10, below right=0.55cm and 0.2cm of main] (llm) {\textbf{llm\_client.py}\\{\scriptsize \texttt{query()}}};

\node[mod, fill=red!8, below=1.35cm of ctx] (val) {\textbf{safety\_validator.py}\\{\scriptsize \texttt{validate()}}};
\node[mod, fill=gray!14, left=0.35cm of val] (cfg) {\textbf{config.py}\\{\scriptsize constants, env}};

\draw[arr] (main) -- (uart);
\draw[arr] (main) -- (ctx);
\draw[arr] (main) -- (llm);
\draw[arr] (main) -- (val);
\draw[arr] (llm.south) -- ++(0,-0.12) -| (cfg.north);
\draw[arr] (val.west) -- (cfg.east);
\draw[arr] (uart.south) -- ++(0,-0.12) -| (cfg.north);

\begin{scope}[on background layer]
  \node[grp, fill=blue!3, draw=blue!35, fit=(main), inner sep=10pt, label={[font=\scriptsize\itshape, text=blue!55]above:Orchestration}] {};
  \node[grp, fill=green!4, draw=green!40, fit=(uart)(ctx)(llm), inner sep=10pt, label={[font=\scriptsize\itshape, text=green!55]above:I/O and reasoning}] {};
  \node[grp, fill=gray!4, draw=gray!45, fit=(val)(cfg), inner sep=10pt, label={[font=\scriptsize\itshape, text=gray!55]below:Safety and configuration}] {};
\end{scope}

\end{tikzpicture}%
}% end resizebox
\caption{Low-level design: Python modules under \texttt{ai\_layer/}}
\label{fig:lld_overview}
\end{figure}

Fig.~\ref{fig:lld_overview} reflects the call pattern described in Section~\ref{sec:modules}: \texttt{main.py} invokes UART bridge, prompt builder, LLM client, and validator; \texttt{config.py} supplies shared limits and environment-backed settings to clients that depend on model id, serial port, and action sets.

%-------------------------------------------------------------------
\section{Use Case Diagram}
\label{sec:usecase}

Figure~\ref{fig:usecase} relates actors to the operational use cases of the deployed robot and cloud service.

\begin{figure}[htbp!]
\centering
\resizebox{0.85\textwidth}{!}{%
\begin{tikzpicture}[
  actor/.style={font=\small},
  uc/.style={draw, thick, ellipse, minimum width=3.6cm, minimum height=0.82cm, align=center, font=\footnotesize, fill=blue!6},
  arr/.style={-{Stealth[length=2mm]}, thick},
  sys/.style={draw, thick, rounded corners=6pt, minimum width=8.5cm, minimum height=8.2cm, fill=gray!3},
  node distance=0.45cm
]

\node[sys, label={[font=\small\bfseries]above:KANDA (field deployment)}] (boundary) {};

\node[uc] at (0, 3.2) (uc1) {Power \& baseline check};
\node[uc] at (0, 2.0) (uc2) {Run intelligence loop (Pi)};
\node[uc] at (0, 0.8) (uc3) {Observe status (OLED / logs)};
\node[uc] at (0, -0.4) (uc4) {Configure port / API key / model};
\node[uc] at (0, -1.6) (uc5) {AUTO\_MODE avoidance only};
\node[uc] at (0, -2.8) (uc6) {AI\_MODE LLM-guided motion};

\node[actor, left=3.1cm of boundary.west] (user) {};
\draw (user) ++(0,0.5) circle(0.25);
\draw (user) ++(0,0.25) -- ++(0,-0.5);
\draw (user) ++(-0.3,0) -- ++(0.6,0);
\draw (user) ++(0,-0.25) -- ++(-0.2,-0.4);
\draw (user) ++(0,-0.25) -- ++(0.2,-0.4);
\node[font=\small\bfseries, below=0.08cm] at ([yshift=-0.62cm]user) {Operator};

\node[actor, right=3.1cm of boundary.east] (dev) {};
\draw (dev) ++(0,0.5) circle(0.25);
\draw (dev) ++(0,0.25) -- ++(0,-0.5);
\draw (dev) ++(-0.3,0) -- ++(0.6,0);
\draw (dev) ++(0,-0.25) -- ++(-0.2,-0.4);
\draw (dev) ++(0,-0.25) -- ++(0.2,-0.4);
\node[font=\small\bfseries, below=0.08cm] at ([yshift=-0.62cm]dev) {Maintainer};

\node[actor, right=2.6cm of boundary.east, yshift=2.3cm] (gemact) {};
\draw (gemact) ++(0,0.45) circle(0.22);
\draw (gemact) ++(0,0.22) -- ++(0,-0.45);
\draw (gemact) ++(-0.25,0) -- ++(0.5,0);
\draw (gemact) ++(0,-0.22) -- ++(-0.18,-0.35);
\draw (gemact) ++(0,-0.22) -- ++(0.18,-0.35);
\node[font=\small\bfseries, below=0.06cm] at ([yshift=-0.55cm]gemact) {Gemini API};

\draw[thick] ([xshift=0.25cm]user.east) -- (uc1.west);
\draw[thick] ([xshift=0.25cm]user.east) -- (uc2.west);
\draw[thick] ([xshift=0.25cm]user.east) -- (uc3.west);
\draw[thick] ([xshift=0.25cm]user.east) -- (uc5.west);
\draw[thick] ([xshift=0.25cm]user.east) -- (uc6.west);

\draw[thick] ([xshift=-0.25cm]dev.west) -- (uc4.east);
\draw[thick] ([xshift=-0.25cm]dev.west) -- (uc2.east);

\draw[thick] ([xshift=-0.22cm]gemact.west) -- (uc2.east);
\draw[thick] ([xshift=-0.22cm]gemact.west) -- (uc6.east);

\end{tikzpicture}%
}% end resizebox
\caption{Use case diagram for KANDA deployment}
\label{fig:usecase}
\end{figure}

The \textit{operator} powers the platform, starts or stops the Pi process, reads the OLED, and relies on obstacle behaviour in AUTO\_MODE or guided behaviour in AI\_MODE. The \textit{maintainer} adjusts serial device paths, credentials, and model identifiers. The \textit{Gemini API} participates whenever the intelligence loop issues HTTPS inference requests.

%-------------------------------------------------------------------
\section{Component Diagram}
\label{sec:component}

Figure~\ref{fig:component} packages software onto the two compute nodes and shows external services.

\begin{figure}[htbp!]
\centering
\resizebox{0.9\textwidth}{!}{%
\begin{tikzpicture}[
  comp/.style={draw, thick, rounded corners=3pt, minimum width=3.8cm, minimum height=1.05cm, align=center, font=\small, fill=white},
  pkg/.style={draw, thick, rounded corners=8pt, minimum width=8.2cm, minimum height=3.4cm},
  arr/.style={-{Stealth[length=2.5mm]}, thick},
  lbl/.style={font=\scriptsize, fill=white, inner sep=1pt},
  node distance=0.75cm
]

\node[comp, fill=cyan!10] (pi) {\textbf{Raspberry Pi 4}\\{\scriptsize Linux + Python 3}};

\node[comp, fill=orange!12, right=2.2cm of pi] (esp) {\textbf{ESP32 DevKit}\\{\scriptsize Arduino core firmware}};

\node[comp, fill=purple!10, below=1.5cm of pi] (net) {\textbf{Network path}\\{\scriptsize outbound HTTPS}};

\node[comp, fill=purple!14, right=0.2cm of net] (gem) {\textbf{Google Gemini}\\{\scriptsize generative API}};

\node[comp, fill=gray!12, below=1.4cm of esp] (bat) {\textbf{Power / mechanics}\\{\scriptsize battery, driver, motors}};

\draw[arr] (pi) -- node[lbl, above] {UART 115200} (esp);
\draw[arr] (pi) -- node[lbl, left] {TLS} (net);
\draw[arr] (net) -- (gem);
\draw[arr] (esp) -- node[lbl, right] {GPIO} (bat);

\begin{scope}[on background layer]
  \node[pkg, fill=cyan!3, draw=cyan!40, fit=(pi), inner sep=14pt, label={[font=\scriptsize\bfseries, text=cyan!60]above left:Companion computer}] {};
  \node[pkg, fill=orange!3, draw=orange!40, fit=(esp)(bat), inner sep=14pt, label={[font=\scriptsize\bfseries, text=orange!60]above right:Embodiment node}] {};
\end{scope}

\end{tikzpicture}%
}% end resizebox
\caption{Component diagram: companion computer, embodiment node, and cloud API}
\label{fig:component}
\end{figure}

The Raspberry Pi component hosts the entire \texttt{ai\_layer/} tree. The ESP32 component hosts sensing, actuation, OLED rendering, and UART framing. Gemini is an external software service accessed only from the Pi.

%-------------------------------------------------------------------
\section{Structure Chart}
\label{sec:structure_chart}

Figure~\ref{fig:structure_chart} shows the static decomposition of the orchestrator and its collaborators.

\begin{figure}[htbp!]
\centering
\resizebox{0.95\textwidth}{!}{%
\begin{tikzpicture}[
  mod/.style={draw, thick, rounded corners=3pt, minimum width=2.85cm, minimum height=0.88cm, align=center, font=\footnotesize},
  top/.style={mod, fill=blue!14, minimum width=5.5cm, minimum height=1.05cm, font=\small\bfseries},
  arr/.style={-{Stealth[length=2mm]}, thick},
  lbl/.style={font=\tiny, fill=white, inner sep=1pt},
  node distance=0.55cm
]

\node[top] (main) {main.py\\\scriptsize process entry, signal handlers};

\node[mod, fill=green!10, below left=0.95cm and 0.15cm of main] (uart) {uart\_bridge.py\\\scriptsize serial I/O};
\node[mod, fill=green!10, below=0.95cm of main] (ctx) {context\_builder.py\\\scriptsize prompts};
\node[mod, fill=green!10, below right=0.95cm and 0.15cm of main] (llm) {llm\_client.py\\\scriptsize Gemini};

\node[mod, fill=red!9, below=1.05cm of ctx] (val) {safety\_validator.py\\\scriptsize validate()};

\node[mod, fill=gray!13, left=0.4cm of val] (cfg) {config.py\\\scriptsize env + limits};

\draw[arr] (main) -- (uart);
\draw[arr] (main) -- (ctx);
\draw[arr] (main) -- (llm);
\draw[arr] (main) -- (val);
\draw[arr] (llm.south) -- ++(0,-0.15) -| (cfg.north);
\draw[arr] (val.west) -- (cfg.east);
\draw[arr] (uart.south) -- ++(0,-0.15) -| (cfg.north);

\end{tikzpicture}%
}% end resizebox
\caption{Structure chart: intelligence layer modules}
\label{fig:structure_chart}
\end{figure}

Entry resolves in \texttt{main.py}: an infinite loop calls \texttt{run\_cycle()}, which sequences UART reads, prompt construction, Gemini inference, validation, and command transmission. Shared constants and limits live in \texttt{config.py}.

%-------------------------------------------------------------------
\section{Control Cycle Activity View}
\label{sec:activity}

Figure~\ref{fig:activity_cycle} refines the Level~2 intelligence pipeline from Chapter~\ref{ch:hld} into the conditional branches implemented in code.

\begin{figure}[htbp!]
\centering
\resizebox{0.72\textwidth}{!}{%
\begin{tikzpicture}[
  flowstep/.style={draw, thick, rounded corners=3pt, minimum width=6.8cm, minimum height=0.62cm, align=center, font=\footnotesize, fill=blue!6},
  dec/.style={draw, thick, diamond, aspect=2, minimum width=3.2cm, inner sep=2pt, align=center, font=\footnotesize, fill=yellow!15},
  arr/.style={-{Stealth[length=2mm]}, thick},
  lbl/.style={font=\tiny, fill=white, inner sep=1pt},
  node distance=0.42cm
]

\node[flowstep] (a) {Read one telemetry line via \texttt{UARTBridge.read\_telemetry()}};
\node[dec, below=0.5cm of a] (b) {Parsed?};
\node[flowstep, below left=0.6cm and 0.1cm of b] (c) {Send safe stop (\texttt{stop}, 0)};
\node[flowstep, below right=0.6cm and 0.1cm of b] (d) {\texttt{build\_prompt(telemetry)}};
\node[flowstep, below=0.45cm of d] (e) {\texttt{llm\_client.query(prompt)}};
\node[flowstep, below=0.45cm of e] (f) {\texttt{safety\_validator.validate(raw)}};
\node[flowstep, below=0.45cm of f] (g) {\texttt{send\_command(action, speed)}};
\node[flowstep, below=0.45cm of g] (h) {Sleep \texttt{LOOP\_INTERVAL\_SEC}; repeat};

\draw[arr] (a) -- (b);
\draw[arr] (b) -- node[lbl, above left] {no} (c);
\draw[arr] (b) -- node[lbl, above right] {yes} (d);
\draw[arr] (d) -- (e);
\draw[arr] (e) -- (f);
\draw[arr] (f) -- (g);
\draw[arr] (g) -- (h);
\draw[arr] (c.south) |- ++(0,-1.8) -| (h.south);

\end{tikzpicture}%
}% end resizebox
\caption{Activity view of one intelligence-layer control cycle (\texttt{run\_cycle})}
\label{fig:activity_cycle}
\end{figure}

Timeouts and API exceptions on the Gemini step route to the same safe-stop policy as missing telemetry before the next interval begins.

%-------------------------------------------------------------------
\section{Module Description}
\label{sec:modules}

Table~\ref{tab:modules} lists Python modules under \texttt{ai\_layer/}. The embodiment artefact is the firmware sketch \texttt{phase2\_obstacle\_avoidance.ino} described in Section~\ref{subsec:firmware}.

\begin{table}[htbp!]
\centering
\caption{Intelligence-layer modules and responsibilities}
\label{tab:modules}
\begin{tabularx}{\textwidth}{@{}>{\raggedright\arraybackslash}p{2.6cm}>{\raggedright\arraybackslash}p{3.4cm}X@{}}
\toprule
\textbf{Module} & \textbf{File} & \textbf{Primary responsibility} \\
\midrule
Orchestrator & \texttt{main.py} & Control loop, signals, shutdown stop, invokes cycle steps \\
Configuration & \texttt{config.py} & Serial port, baud, model id, timing, \texttt{VALID\_ACTIONS}, speed bounds \\
UART bridge & \texttt{uart\_bridge.py} & Open serial, parse telemetry lines, send JSON commands \\
Context builder & \texttt{context\_builder.py} & Static hardware description + live sensor section + Phase~4 stubs \\
LLM client & \texttt{llm\_client.py} & Configure Gemini, generate content, extract JSON from response text \\
Safety validator & \texttt{safety\_validator.py} & Type and membership checks; clamp speed; fallback stop dict \\
\bottomrule
\end{tabularx}
\end{table}

%-------------------------------------------------------------------
\subsection[Configuration module]{Configuration module (\texttt{config.py})}
\label{subsec:config}

\textbf{Responsibility:} Centralise tunable parameters and environment-backed secrets so other modules avoid scattering literals.

\textbf{Key items:} \texttt{SERIAL\_PORT} (default \texttt{/dev/ttyS0}), \texttt{BAUD\_RATE} (115200), \texttt{GEMINI\_MODEL} (default \texttt{gemini-2.5-flash}), \texttt{GEMINI\_API\_KEY} from environment, \texttt{VALID\_ACTIONS} as a Python set of seven string verbs, \texttt{SPEED\_MIN}/\texttt{SPEED\_MAX} (0--255), \texttt{LOOP\_INTERVAL\_SEC}, \texttt{LLM\_TIMEOUT\_SEC}, and \texttt{DEFAULT\_SPEED} when the model omits speed.

%-------------------------------------------------------------------
\subsection[UART bridge module]{UART bridge module (\texttt{uart\_bridge.py})}
\label{subsec:uart}

\textbf{Responsibility:} Encapsulate pyserial lifecycle and the ASCII protocols agreed with firmware.

\textbf{Interface:} Class \texttt{UARTBridge} exposes \texttt{connect()}, \texttt{disconnect()}, context-manager support, \texttt{send\_command(action, speed)} (serialises a one-line JSON object with \texttt{action} and \texttt{speed} keys), and \texttt{read\_telemetry()} which blocks up to \texttt{SERIAL\_TIMEOUT}, decodes UTF-8, and parses lines of the form \texttt{F:<cm> L:<cm> R:<cm> -> <ACTION>} into a dictionary with floating-point ranges and a string action tag.

\textbf{Notes:} A short sleep after opening the port accommodates ESP32 reset behaviour.

%-------------------------------------------------------------------
\subsection[Context builder module]{Context builder module (\texttt{context\_builder.py})}
\label{subsec:ctx}

\textbf{Responsibility:} Concatenate a static hardware-and-policy preamble with formatted front/left/right readings so Gemini receives explicit affordances and safety hints on every call~\cite{ahn2022saycan}.

\textbf{Interface:} \texttt{build\_prompt(telemetry, user\_speech=None, image\_b64=None)} returns one prompt string. Optional speech and image parameters are reserved for Phase~4 multimodal extensions; the LLM client already accepts \texttt{image\_b64} for future attachment.

%-------------------------------------------------------------------
\subsection[LLM client module]{LLM client module (\texttt{llm\_client.py})}
\label{subsec:llm}

\textbf{Responsibility:} Initialise the Google Generative AI client once per process, invoke \texttt{GenerativeModel.generate\_content} with low temperature (0.2) and short \texttt{max\_output\_tokens} (64), and recover a JSON object from arbitrary response text.

\textbf{Parsing strategy:} Prefer fenced markdown JSON if present; otherwise extract the first balanced brace segment and pass it to \texttt{json.loads}. Failures surface as \texttt{ValueError} for the orchestrator to translate into safe stop.

%-------------------------------------------------------------------
\subsection[Safety validator module]{Safety validator module (\texttt{safety\_validator.py})}
\label{subsec:val}

\textbf{Responsibility:} Guarantee that only whitelisted actions and bounded speeds reach \texttt{send\_command}; never raise---always return a usable dictionary~\cite{huang2023grounded}.

\textbf{Rules:} Reject non-dicts and unknown actions with fallback \texttt{action}\,=\,\texttt{stop} and \texttt{speed}\,=\,0; coerce \texttt{speed} with \texttt{int()}; clamp to $[0,255]$ with logging when out of range.

%-------------------------------------------------------------------
\subsection[Orchestrator]{Orchestrator (\texttt{main.py})}
\label{subsec:main}

\textbf{Responsibility:} Own the process lifecycle: configure logging, register SIGINT/SIGTERM handlers that send stop, open UART, loop with \texttt{time.sleep(LOOP\_INTERVAL\_SEC)}, catch unexpected exceptions per cycle and issue stop, and disconnect cleanly.

\textbf{Cycle ordering:} Matches Fig.~\ref{fig:activity_cycle}: read telemetry; if missing, stop and skip LLM; else build prompt, query Gemini (errors treated like invalid output), validate, transmit.

%-------------------------------------------------------------------
\subsection[Embodiment firmware]{Embodiment firmware (\texttt{phase2\_obstacle\_avoidance.ino})}
\label{subsec:firmware}

\textbf{Responsibility:} Run the real-time sense--act loop on the ESP32: trigger/echo ultrasonic measurements with \texttt{pulseIn}, map distances through motor primitives (\texttt{forward}, pivots, etc.) via TB6612FNG direction pins and LEDC PWM, refresh the SSD1306 via I\textsuperscript{2}C, and exchange newline-terminated messages over \texttt{Serial}.

\textbf{Modes:} \texttt{AI\_MODE} activates when a valid JSON line arrives (\texttt{ArduinoJson} parse); \texttt{AUTO\_MODE} executes distance-threshold rules when no JSON is seen for \texttt{AI\_TIMEOUT\_MS} (5000\,ms). Constants such as \texttt{OBSTACLE\_DIST} and \texttt{WALL\_DIST} tune avoidance aggressiveness independently of the LLM policy on the Pi~\cite{brooks1986robust}.

%-------------------------------------------------------------------
\section{Summary}
\label{sec:ch5summary}

This chapter detailed KANDA's software structure. Figures~\ref{fig:hld_overview}--\ref{fig:component} positioned the Pi intelligence layer, ESP32 embodiment node, UART link, and Gemini service; Figures~\ref{fig:structure_chart} and~\ref{fig:activity_cycle} traced module ownership and one control cycle. Table~\ref{tab:modules} and Sections~\ref{subsec:config}--\ref{subsec:firmware} documented configuration, serial bridging, prompting, Gemini access, validation, orchestration, and firmware responsibilities consistent with Chapters~\ref{ch:srs} and~\ref{ch:hld}. Implementation detail appears in Chapter~\ref{ch:impl}; formal test procedures and requirement traceability appear in Chapter~\ref{ch:testing}.
